\documentclass{article}
\usepackage{amsmath,amssymb}
\usepackage[vocab={majorization}]{leanify}
\newtheorem{lemma}{Lemma}
\begin{document}

\section{Notation}

We write $L_k(x)$ for the $k$th largest coordinate of $x$.
We write $S_k(x)$ for the sum of the $k$ largest coordinates of $x$.
We write $\Pi(x,t)$ for the positive-part sum of $x$ at $t$.

\section{Majorization}

\begin{lemma}\label{reflexive}
Let $n$ be a positive integer and let $x$ be a vector in $\mathbb{R}^n$. Then $x$ is majorized by $x$, and $S_0(x)=0$.
\end{lemma}

\begin{lemma}\label{weak}
Let $n$ be a positive integer, and let $x$ and $y$ be vectors in $\mathbb{R}^n$. If $x$ is majorized by $y$, then $x$ is weakly submajorized by $y$.
\end{lemma}

\begin{lemma}\label{order}
Let $x$ be a vector in $\mathbb{R}^2$. Then $L_1(x)\ge L_2(x)$, and $\Pi(x,t)\ge0$ for every real number $t$.
\end{lemma}

\end{document}
